\scanpage{119}\mspage{64}
On voit par là qu'on a des isom. canoniques
\[
 \unsure{\rho|p}\,i\simeq\mathrm{id},\qquad \sigma j\simeq\mathrm{id}
\]
dans la catégorie des endofoncteurs de $(\mathrm{Cat}_{/S})W_S^{-1}$ $(\approx\mathrm{HOT}^{\unsure{\mathrm{lc}|\ell\mathrm{c}}}_W(S))$, \deleted{\unsure{puisque}} isomorphismes qui s'explicitent ici par \deleted{\unsure{l'isom} \unsure{homotopique|homotopie}} \inserted{le fait} que $\rho i,\sigma j$ sont déduits du foncteur $\unsure{C|\mathcal C}$ par passage \unsure{aux localisés|au localisé} \inserted{$W_S^{-1}$} (donc en fait $\rho i=\sigma j$ !), et qu'on a une \deleted{\unsure{isom}} flèche fonctorielle en $X$
\[
 X\xrightarrow{\unsure{i_x|i_X}}C(X),
\]
qui est dans $W$ pour tt $X$. Ainsi, $\rho$ et $\sigma$ sont \inserted{à isom près} des \uline{rétractions} des \deleted{\unsure{C}} catégories $\mathrm{HOT}_W(S)$, $\mathrm{HOT}_W^\circ(S)$ sur la catégorie $\mathrm{HOT}_W^{\mathrm{lc}}(S)$. \deleted{Celle-ci} De \deleted{est \unsure{équivale}} plus, ces rétractions sont \deleted{\unsure{des}} \deleted{\unsure{c. à dire|c'est à dire}} ce qu'on appelle des foncteurs de localisation (comme on voit sur le diagramme de la page 62). (On dit que $C\xrightarrow{F}C'$ est un \uline{foncteur de localisation} si, désignant par $\Sigma\subset\mathrm{Fl}\,C$ l'ens des flèches de $C$ transformées en des iso par $F$, le foncteur
\[
 C\Sigma^{-1}\xrightarrow{\overline F}C'
\]
qui factorise $F$ est une équivalence de catégories.)
% archive sheet number: 118
